\labsection{6}{Quantify uncertainty with bootstrap and permutation logic}{sec:lab06}

\begin{labproject}{Sampling and inference: estimate, quantify, challenge}
\textbf{Coverage.} Chapter~\ref{ch:inference}.\\
\textbf{Goal.} You should be able to define an estimand, construct a bootstrap interval, build a permutation reference distribution, report effect size, and separate statistical evidence from study-design limitations.\\
\textbf{Dataset.} \texttt{student\_performance\_dataset.csv}.\\
\textbf{Scenario.} A school analyst observes a grade difference between students who report intending higher education and those who do not. Your task is to quantify the difference, assess sampling uncertainty, and prevent statistical significance from being mistaken for a causal conclusion.

\textbf{Core workflow.}
\begin{enumerate}
  \item \textbf{Define the estimand.} Let the target be the difference in mean final grade \texttt{G3} between \texttt{higher=yes} and \texttt{higher=no}. Report group sizes and descriptive summaries before inference.
  \item \textbf{Estimate the observed effect.} Compute the mean difference and a standardised difference. Interpret direction and practical size in the units of the problem.
  \item \textbf{Bootstrap uncertainty.} Resample within groups, generate at least 5,000 bootstrap estimates, and construct a 95\% percentile interval for the mean difference.
  \item \textbf{Construct a null model.} Permute the group labels while keeping grades fixed. Use the resulting distribution to estimate a two-sided p-value for no group association.
  \item \textbf{Compare inferential quantities.} Explain what the point estimate, confidence interval, standardised effect, and p-value each contribute. Explicitly state what none of them establishes.
  \item \textbf{Audit design limitations.} Identify at least three available variables that may confound the association, such as prior grades, failures, study time, family education, or absences.
  \item \textbf{Write a defensible conclusion.} Separate the statistical conclusion from the causal conclusion in two distinct sentences.
\end{enumerate}

\begin{tryit}
\textbf{Extension.} Repeat the comparison after stratifying or adjusting descriptively for one plausible confounder. Do not claim that this single adjustment removes confounding; use it as a sensitivity check on the original result.
\end{tryit}

\tcbbreak
\lstinputlisting[style=utpr]{all_experiments/lab06_inference.R}
\begin{verification}
\textbf{Required evidence.} Submit the observed effect, 95\% interval, permutation p-value, one bootstrap distribution graphic, and a four-sentence interpretation covering magnitude, uncertainty, null evidence, and design limitation.

\textbf{Professional standard.} Report uncertainty and effect size together. A small p-value without a meaningful effect estimate is not a complete scientific result.
\end{verification}
\end{labproject}
